\cotefile{157-5}
\scanpage{633}
\mspage{21}

à vérifier seulement

\[
 v^{*}v_{!}(TC)\subset TC,\tag*{(2.3.7)}
\]

or on a, pour $X=(X_i)_{i\in I_0}\in M(I_0)$

\[
\begin{aligned}
 v_!(X)(i)&\simeq\coprod_{j\in I_0}\unsure{v|u}_{j\,!}(X_j)(i)\\
 &\simeq\coprod_{j\in I_0}\underbrace{X_j^{(\operatorname{Hom}(i,j))}}_{\substack{\text{somme directe}\\\text{de }\operatorname{Hom}(i,j)\text{ copies}\\\text{de }X_j}},
\end{aligned}
\]

et la relation (2.3.7) résulte du fait que
$TC$, \inserted{\unsure{étant|il est} $Q$-clos à g.} est stable par petites sommes directes.

\uline{Prouvons $2^\circ$}. Les formules (2.3.2) sont
essentiellement tautologiques, les deux dernières
\unsure{résultent|résultant} de la troisième et de la deuxième,
en vertu de 1.5 appliqué au couple de
foncteurs adjoints $(u_!,u^*)$. La relation
$u_!(C\langle I\rangle)\subset C\langle I'\rangle$, et le fait que
$u_!$ transforme objet initial en objet initial,
impliquent (2.3.3). \unsure{Que} $u_!$ restreint à
$M(I)_c$ soit compatible aux localiseurs \unsure{provient}
de $u_!(TC\langle I\rangle)\subset TC\langle I'\rangle$, et du
fait que \unsure{tte|toute} flèche $f$ dans $M(I)_c$ qui est
dans $W(I)$, se décompose en $f=pi$, avec
$i\in TC$ et $p$ inverse à gauche d'une $j\in TC\langle I\rangle$.
Ceci \deleted{\unsure{dit}}\inserted{vu}, \deleted{\unsure{les}}\inserted{\unsure{l'une des|l'une}} deux premières formules
(2.3.2), montre qu'on est sous les conditions
% archive sheet number: 632
